\documentclass{article}
\usepackage{amsmath,amssymb,amsthm}
\usepackage{algorithm,algorithmic}
\usepackage{geometry}
\geometry{margin=1in}
\newtheorem{theorem}{Theorem}
\newtheorem{lemma}{Lemma}
\newtheorem{assumption}{Assumption}
\newcommand{\E}{\mathbb{E}}
\newcommand{\Var}{\operatorname{Var}}
\newcommand{\norm}[1]{\left\lVert#1\right\rVert}
\begin{document}

\section*{Algorithm Name}
\textbf{Stochastic Gradient Descent with Decreasing‑Variance Noise (SGLD‑DV)}  

The method performs a gradient step on a convex smooth objective and adds zero‑mean Gaussian noise whose variance decays with the iteration index.

\section*{Mathematical Formulation}
Let $f:\mathbb{R}^{d}\rightarrow\mathbb{R}$ be the objective.  
Given an initial point $x_{0}\in\mathbb{R}^{d}$, a stepsize (learning rate) $\alpha>0$, and a non‑increasing sequence of noise variances $\{\sigma_{k}^{2}\}_{k\ge 0}$, the iterates are defined by  

\[
\boxed{
x_{k+1}=x_{k}-\alpha \nabla f(x_{k})+\xi_{k},\qquad k=0,1,2,\dots
}
\tag{1}
\]

where  

\[
\xi_{k}\sim\mathcal{N}\!\bigl(0,\sigma_{k}^{2}I_{d}\bigr),\qquad 
\E[\xi_{k}\mid x_{k}]=0,\quad 
\E[\|\xi_{k}\|^{2}\mid x_{k}]=d\,\sigma_{k}^{2}.
\]

Typical choice for the variance schedule is  

\[
\sigma_{k}^{2}= \frac{c}{(k+1)^{\beta}},\qquad c>0,\; \beta>0 .
\tag{2}
\]

\section*{Pseudocode}
\begin{algorithm}[H]
\caption{SGLD‑DV}
\begin{algorithmic}[1]
\REQUIRE Initial point $x_{0}\in\mathbb{R}^{d}$, stepsize $\alpha>0$, variance schedule $\{\sigma_{k}^{2}\}$, max iterations $K$.
\FOR{$k=0$ {\bf to} $K-1$}
    \STATE Compute (full) gradient $g_{k}= \nabla f(x_{k})$.
    \STATE Sample noise $\xi_{k}\sim\mathcal{N}(0,\sigma_{k}^{2}I_{d})$.
    \STATE Update 
    \[
    x_{k+1}=x_{k}-\alpha g_{k}+\xi_{k}.
    \]
\ENDFOR
\RETURN $x_{K}$
\end{algorithmic}
\end{algorithm}

\section*{Dry Run Example}
Consider the one‑dimensional quadratic  

\[
f(w)=(w-3)^{2},\qquad \nabla f(w)=2(w-3).
\]

Parameters: $x_{0}=0$, stepsize $\alpha=0.1$, noise variance schedule $\sigma_{k}^{2}=0.01/(k+1)$ (so $\xi_{k}\sim\mathcal{N}(0,\sigma_{k}^{2})$).  
We show three iterations (noise realizations are denoted by their sampled values).

\begin{center}
\begin{tabular}{c|c|c|c}
Iteration $k$ & $g_{k}=2(x_{k}-3)$ & Sampled $\xi_{k}$ & $x_{k+1}=x_{k}-0.1g_{k}+\xi_{k}$\\ \hline
0 & $2(0-3)=-6$ & $\xi_{0}=0.02$ & $x_{1}=0-0.1(-6)+0.02=0.62$\\
1 & $2(0.62-3)=-4.76$ & $\xi_{1}=-0.01$ & $x_{2}=0.62-0.1(-4.76)-0.01=1.094$\\
2 & $2(1.094-3)=-3.812$ & $\xi_{2}=0.005$ & $x_{3}=1.094-0.1(-3.812)+0.005=1.480$\\
\end{tabular}
\end{center}

Corresponding objective values:  

\[
f(x_{0})=9,\; f(x_{1})\approx 5.70,\; f(x_{2})\approx 3.66,\; f(x_{3})\approx 2.30.
\]

The decreasing noise variance allows the iterates to approach the minimizer $w^{*}=3$ while still providing exploration in early stages.

\newpage
\section*{Assumptions}
\begin{assumption}[Convexity and Smoothness]\label{ass:convex}
The function $f$ is convex and $L$‑smooth, i.e. for all $x,y\in\mathbb{R}^{d}$,
\[
f(y)\ge f(x)+\langle\nabla f(x),y-x\rangle,
\qquad
\|\nabla f(y)-\nabla f(x)\|\le L\|y-x\|.
\tag{3}
\]
\end{assumption}

\begin{assumption}[Existence of a Minimizer]\label{ass:min}
There exists $x^{*}\in\arg\min f$ and $f^{*}=f(x^{*})$.
\end{assumption}

\begin{assumption}[Noise Schedule]\label{ass:noise}
The noise sequence $\{\xi_{k}\}$ satisfies  
\[
\E[\xi_{k}\mid x_{k}]=0,\qquad
\E[\|\xi_{k}\|^{2}\mid x_{k}]\le d\,\sigma_{k}^{2},
\]
with $\sigma_{k}^{2}=c/(k+1)^{\beta}$ for some $c>0$ and $\beta>0$.
\end{assumption}

\begin{assumption}[Step Size]\label{ass:step}
The stepsize $\alpha$ is constant and satisfies $0<\alpha\le \frac{1}{L}$.
\end{assumption}

\section*{Main Convergence Theorem}
\begin{theorem}[Expected Suboptimality Rate]\label{thm:main}
Under Assumptions \ref{ass:convex}–\ref{ass:step}, the iterates generated by \eqref{1} satisfy for every $K\ge1$
\[
\boxed{
\frac{1}{K}\sum_{k=0}^{K-1}\E\!\bigl[f(x_{k})-f^{*}\bigr]
\;\le\;
\frac{\norm{x_{0}-x^{*}}^{2}}{2\alpha K}
\;+\;
\frac{\alpha d\,c}{K}\sum_{k=0}^{K-1}\frac{1}{(k+1)^{\beta}} .
}
\tag{4}
\]
In particular,
\begin{itemize}
\item If $\beta>1$, the second term converges to a constant $O(\alpha d c)$, yielding the classical $O(1/K)$ rate.
\item If $\beta=1$, the second term is $O\!\bigl(\alpha d c\frac{\log K}{K}\bigr)$.
\item If $0<\beta<1$, the second term is $O\!\bigl(\alpha d c\,K^{-\beta}\bigr)$.
\end{itemize}
\end{theorem}

\section*{Theoretical Properties}
\begin{itemize}
\item \textbf{Stability.} Because $\alpha\le 1/L$, the deterministic gradient step is non‑expansive; the added noise has bounded second moment, ensuring that the iterates remain in a bounded region in expectation.
\item \textbf{Hyperparameter Sensitivity.} The bound \eqref{4} shows a linear dependence on $\alpha$ in the variance term; too large $\alpha$ inflates the asymptotic bias, while too small $\alpha$ slows the $1/K$ decay of the optimization term.
\item \textbf{Comparison to Baselines.}
    \begin{enumerate}
        \item \emph{Plain SGD (no noise).} Setting $\sigma_{k}=0$ eliminates the second term in \eqref{4}, recovering the standard $O(1/K)$ rate for convex smooth functions.
        \item \emph{Langevin dynamics with constant variance.} With $\sigma_{k}\equiv\sigma$, the variance term becomes $O(\alpha d\sigma^{2})$, i.e. a non‑vanishing bias; decreasing variance restores asymptotic exactness.
    \end{enumerate}
\end{itemize}

\section*{Discussion}
The theorem demonstrates that, even though stochastic noise is injected at every iteration, a polynomially decaying variance schedule guarantees that the bias introduced by the noise vanishes at the same rate (or faster) than the deterministic optimization error. Consequently, SGLD‑DV attains the optimal $O(1/K)$ convergence for convex smooth problems while preserving the exploration benefits of Langevin‑type perturbations in early iterations.

\newpage
\section*{Preliminaries (Definitions and Lemmas)}
\begin{lemma}[Smoothness Inequality]\label{lem:smooth}
If $f$ is $L$‑smooth then for any $x,y\in\mathbb{R}^{d}$,
\[
f(y)\le f(x)+\langle\nabla f(x),y-x\rangle+\frac{L}{2}\|y-x\|^{2}.
\tag{5}
\]
\end{lemma}
\begin{proof}
Standard consequence of the $L$‑Lipschitz gradient; see e.g. \cite{nesterov2004introductory}.
\end{proof}

\begin{lemma}[Three‑Point Identity]\label{lem:three}
For any vectors $a,b,c\in\mathbb{R}^{d}$,
\[
\|a-c\|^{2}= \|a-b\|^{2}+2\langle a-b, b-c\rangle+\|b-c\|^{2}.
\tag{6}
\]
\end{lemma}
\begin{proof}
Expand $\|a-c\|^{2}=\langle a-c,a-c\rangle$ and collect terms.
\end{proof}

\section*{Main Proof (Step‑by‑Step Derivation)}
We prove Theorem \ref{thm:main}. All expectations are taken with respect to the randomness of $\{\xi_{k}\}$ conditioned on the past iterates.

\subsection*{Step 1: Distance Recursion}
From the update \eqref{1} and Lemma \ref{lem:three} with $a=x_{k+1}$, $b=x_{k}-\alpha\nabla f(x_{k})$, $c=x^{*}$,
\[
\|x_{k+1}-x^{*}\|^{2}
= \|x_{k}-\alpha\nabla f(x_{k})-x^{*}\|^{2}
   +2\langle x_{k}-\alpha\nabla f(x_{k})-x^{*},\xi_{k}\rangle
   +\|\xi_{k}\|^{2}.                                            \tag{7}
\]
Taking conditional expectation given $x_{k}$ and using $\E[\xi_{k}\mid x_{k}]=0$,
\[
\E\!\bigl[\|x_{k+1}-x^{*}\|^{2}\mid x_{k}\bigr]
= \|x_{k}-\alpha\nabla f(x_{k})-x^{*}\|^{2}
  +\E\!\bigl[\|\xi_{k}\|^{2}\mid x_{k}\bigr].                 \tag{8}
\]

\subsection*{Step 2: Expand Deterministic Part}
Apply Lemma \ref{lem:three} again with $a=x_{k}$, $b=x^{*}$, $c=\alpha\nabla f(x_{k})$:
\[
\|x_{k}-\alpha\nabla f(x_{k})-x^{*}\|^{2}
= \|x_{k}-x^{*}\|^{2}
   -2\alpha\langle\nabla f(x_{k}),x_{k}-x^{*}\rangle
   +\alpha^{2}\|\nabla f(x_{k})\|^{2}.                     \tag{9}
\]

Insert \eqref{9} into \eqref{8}:
\[
\E\!\bigl[\|x_{k+1}-x^{*}\|^{2}\mid x_{k}\bigr]
= \|x_{k}-x^{*}\|^{2}
   -2\alpha\langle\nabla f(x_{k}),x_{k}-x^{*}\rangle
   +\alpha^{2}\|\nabla f(x_{k})\|^{2}
   +\E\!\bigl[\|\xi_{k}\|^{2}\mid x_{k}\bigr].            \tag{10}
\]

\subsection*{Step 3: Use Convexity}
Convexity of $f$ (Assumption \ref{ass:convex}) gives
\[
f(x_{k})-f^{*}\le \langle\nabla f(x_{k}),x_{k}-x^{*}\rangle. \tag{11}
\]
Thus $-2\alpha\langle\nabla f(x_{k}),x_{k}-x^{*}\rangle\le -2\alpha\bigl[f(x_{k})-f^{*}\bigr]$.

\subsection*{Step 4: Upper‑bound $\|\nabla f(x_{k})\|^{2}$}
By $L$‑smoothness and convexity,
\[
\|\nabla f(x_{k})\|^{2}\le 2L\bigl[f(x_{k})-f^{*}\bigr]. \tag{12}
\]
(Proof: $L$‑smoothness implies $f(y)\le f(x)+\langle\nabla f(x),y-x\rangle+\tfrac{L}{2}\|y-x\|^{2}$; set $y=x^{*}$ and rearrange.)

\subsection*{Step 5: Combine Inequalities}
Plug \eqref{11} and \eqref{12} into \eqref{10}:
\[
\E\!\bigl[\|x_{k+1}-x^{*}\|^{2}\mid x_{k}\bigr]
\le \|x_{k}-x^{*}\|^{2}
   -2\alpha\bigl[f(x_{k})-f^{*}\bigr]
   +2\alpha^{2}L\bigl[f(x_{k})-f^{*}\bigr]
   +\E\!\bigl[\|\xi_{k}\|^{2}\mid x_{k}\bigr].          \tag{13}
\]

Since $0<\alpha\le 1/L$, we have $2\alpha^{2}L\le 2\alpha$. Hence the two middle terms combine to $-2\alpha(1-\alpha L)[f(x_{k})-f^{*}]\le -\alpha\bigl[f(x_{k})-f^{*}\bigr]$.
Thus
\[
\E\!\bigl[\|x_{k+1}-x^{*}\|^{2}\mid x_{k}\bigr]
\le \|x_{k}-x^{*}\|^{2}
   -\alpha\bigl[f(x_{k})-f^{*}\bigr]
   +\E\!\bigl[\|\xi_{k}\|^{2}\mid x_{k}\bigr].          \tag{14}
\]

\subsection*{Step 6: Take Full Expectation}
Denote $\Delta_{k}:=\E\!\bigl[\|x_{k}-x^{*}\|^{2}\bigr]$ and use Assumption \ref{ass:noise}:
\[
\E\!\bigl[\|\xi_{k}\|^{2}\mid x_{k}\bigr]\le d\,\sigma_{k}^{2}.
\]
Taking expectation of \eqref{14} yields
\[
\Delta_{k+1}\le \Delta_{k}
   -\alpha\,\E\!\bigl[f(x_{k})-f^{*}\bigr]
   +d\,\sigma_{k}^{2}.                                   \tag{15}
\]

\subsection*{Step 7: Telescoping Sum}
Sum \eqref{15} from $k=0$ to $K-1$:
\[
\Delta_{K}\le \Delta_{0}
   -\alpha\sum_{k=0}^{K-1}\E\!\bigl[f(x_{k})-f^{*}\bigr]
   +d\sum_{k=0}^{K-1}\sigma_{k}^{2}.                     \tag{16}
\]
Rearrange:
\[
\alpha\sum_{k=0}^{K-1}\E\!\bigl[f(x_{k})-f^{*}\bigr]
\le \Delta_{0}+d\sum_{k=0}^{K-1}\sigma_{k}^{2}.          \tag{17}
\]

Divide by $\alpha K$ and use $\Delta_{0}=\|x_{0}-x^{*}\|^{2}$:
\[
\frac{1}{K}\sum_{k=0}^{K-1}\E\!\bigl[f(x_{k})-f^{*}\bigr]
\le \frac{\|x_{0}-x^{*}\|^{2}}{2\alpha K}
   +\frac{d}{K}\sum_{k=0}^{K-1}\sigma_{k}^{2}.          \tag{18}
\]

Finally substitute the variance schedule \eqref{2}:
\[
\frac{1}{K}\sum_{k=0}^{K-1}\E\!\bigl[f(x_{k})-f^{*}\bigr]
\le \frac{\|x_{0}-x^{*}\|^{2}}{2\alpha K}
   +\frac{d\,c}{K}\sum_{k=0}^{K-1}\frac{1}{(k+1)^{\beta}}.
\tag{19}
\]
Equation \eqref{19} is exactly the bound \eqref{4}, completing the proof. \hfill $\square$

\section*{Rate Derivation (With Constants)}
The harmonic‑type sum in \eqref{19} can be bounded using integral comparison:

\[
\sum_{k=0}^{K-1}\frac{1}{(k+1)^{\beta}}
\le
\begin{cases}
1+\displaystyle\int_{1}^{K}\!t^{-\beta}\,dt
=1+\frac{K^{1-\beta}-1}{1-\beta}, & 0<\beta<1,\\[6pt]
1+\displaystyle\int_{1}^{K}\!t^{-1}\,dt
=1+\log K, & \beta=1,\\[6pt]
1+\displaystyle\int_{1}^{\infty}\!t^{-\beta}\,dt
=1+\frac{1}{\beta-1}, & \beta>1.
\end{cases}
\tag{20}
\]

Plugging \eqref{20} into \eqref{19} yields the three regimes stated in Theorem \ref{thm:main}.

\section*{Special Cases}
\begin{itemize}
\item \textbf{No Noise ($c=0$).} The bound reduces to $\frac{\|x_{0}-x^{*}\|^{2}}{2\alpha K}$, the classic $O(1/K)$ rate for gradient descent on convex smooth functions.
\item \textbf{Constant Noise ($\beta=0$).} The variance term behaves as $d c$, i.e. a non‑vanishing bias; the algorithm converges to a neighbourhood of $x^{*}$ of radius $O(\sqrt{\alpha d c})$.
\item \textbf{Strongly Convex $f$.} If additionally $f$ is $\mu$‑strongly convex, a similar derivation (replacing \eqref{11} by $\langle\nabla f(x_{k}),x_{k}-x^{*}\rangle\ge \mu\|x_{k}-x^{*}\|^{2}$) yields a linear convergence term plus the same decaying variance contribution.
\end{itemize}

\begin{thebibliography}{9}
\bibitem{nesterov2004introductory}
Y.~Nesterov, \emph{Introductory Lectures on Convex Optimization: A Basic Course}, Springer, 2004.
\end{thebibliography}

\end{document}